\chapter{Rester à sa place}
\label{chap:rester-a-sa-place}

Les visites du soir auraient probablement pu continuer encore quelque temps si
elles étaient restées aussi discrètes que Khaer Ghal l'imaginait. Dans un
village de cette taille, cependant, peu de choses demeuraient longtemps
secrètes. On avait fini par remarquer que le fils du chef ne se contentait plus
d'apporter les rations aux prisonniers et qu'il revenait parfois près de leur
enclos lorsque ses tâches étaient terminées.

Son père l'apprit.

Il ne lui en parla pas.

Pendant deux jours, Khaer Ghal continua donc à vivre comme si rien n'avait
changé. Il s'entraîna, participa aux tâches qu'on lui confiait et apporta comme
d'habitude la nourriture aux prisonniers. Le soir, il retrouva Kamatsu et
écouta ses histoires, ignorant que son père savait désormais exactement où il
passait une partie de son temps.

Ce fut seulement le troisième jour qu'il apprit qu'un nouveau combat avait été
annoncé dans l'arène.

La nouvelle n'avait en elle-même rien d'exceptionnel. Les affrontements y
étaient suffisamment fréquents pour que les habitants du village n'y prêtent
guère attention avant de connaître les combattants ou les créatures choisis.
Cette fois, trois prisonniers devaient participer.

Khaer Ghal demanda lesquels.

La réponse lui fit oublier tout le reste.

Kamatsu, sa mère et son père.

Tous les trois seraient conduits ensemble dans la fosse et affronteraient deux
créatures capturées quelques jours auparavant dans les bois.

Khaer Ghal les connaissait. Il en avait déjà observé plusieurs et avait appris
à reconnaître leur manière de chasser : des prédateurs rapides, capables de
bondir brutalement sur une courte distance et dont la mâchoire constituait
l'arme principale. Pour un guerrier entraîné, l'affrontement aurait été
dangereux sans être exceptionnel. Pour trois humains qui n'avaient jamais
appris à combattre ensemble et auxquels on ne donnerait que quelques armes
rudimentaires, les chances étaient très différentes.

Il pensa à Kamatsu lui demandant pourquoi quelqu'un accepterait d'entrer dans
l'arène s'il risquait d'y mourir.

Cette fois, personne ne lui avait demandé son avis.

Khaer Ghal partit trouver son père.

\scenebreak

Le chef du village l'entendit arriver mais ne leva pas immédiatement les yeux.
Khaer Ghal attendit quelques secondes devant lui avant de parler.

— Pourquoi eux ?

Son père finit par le regarder.

— Pourquoi qui ?

— Les trois humains. Ceux qui vont dans l'arène.

— Ce sont des prisonniers.

Quelques semaines auparavant, cette réponse aurait probablement suffi. Les
prisonniers travaillaient, certains combattaient dans l'arène et les décisions
appartenaient aux adultes qui dirigeaient le village. Khaer Ghal avait grandi
avec ces règles et n'avait jamais réellement imaginé qu'elles puissent être
différentes.

Mais Kamatsu avait posé trop de questions.

Et certaines continuaient de revenir.

— Il y en a d'autres.

Son père l'observa plus attentivement.

— Oui.

— Alors pourquoi eux ?

Un silence s'installa entre eux.

— Parce que je l'ai décidé.

Khaer Ghal sentit ses mâchoires se contracter. Il connaissait ce ton. Lorsque
son père parlait ainsi, la discussion était terminée.

Elle aurait dû l'être.

— Choisis-en d'autres.

Le regard du chef se durcit aussitôt.

— Fais attention à la manière dont tu t'adresses à moi.

Khaer Ghal baissa légèrement les yeux par réflexe. Il avait appris très jeune
ce que signifiait ce changement dans la voix de son père, mais cette fois il
ne partit pas.

— Il y en a d'autres, répéta-t-il. Prends-les.

— Pourquoi ?

Khaer Ghal hésita.

Il aurait pu chercher une raison acceptable. Trois prisonniers plus forts
offriraient un meilleur combat. Des adultes seraient plus intéressants à
regarder. Il pouvait probablement trouver quelque chose qui ressemblerait à
un argument.

Mais son père saurait qu'il mentait.

— J'aime bien parler avec le garçon.

Le chef ne répondit pas.

— Il s'appelle Kamatsu.

Quelque chose passa dans son regard.

— Je sais.

Khaer Ghal releva brusquement la tête.

Son père savait. Il savait qu'il revenait près de l'enclos. Il savait qu'il parlait à Kamatsu. Il savait même son nom.

Et soudain le choix de cette famille ne ressemblait plus du tout à un hasard.

— C'est pour ça ?

— Oui.

La réponse fut immédiate.

Khaer Ghal resta interdit quelques secondes.

— Tu avais reçu un ordre, reprit son père. Tu devais nourrir les prisonniers et
ne pas leur adresser la parole. Pourtant, tu t'attardes auprès d'eux. Tu retournes les voir
lorsque personne ne te l'a demandé. Tu leur apportes même ta propre nourriture.

Khaer Ghal ne répondit rien.

— Ce sont des prisonniers. Pas tes compagnons. Pas des membres du clan. Tu n'as
pas à t'attacher à eux.

— Kamatsu n'a rien fait.

— Justement.

Son père se leva.

Khaer Ghal dut relever davantage la tête pour continuer à soutenir son regard.

— Tu dois apprendre qu'un guerrier ne laisse pas ce genre de faiblesse décider
à sa place. Aujourd'hui, tu me demandes de revenir sur une décision pour un
prisonnier. Demain, que feras-tu ? Refuseras-tu de combattre parce que ton
adversaire te plaît ? Laisseras-tu partir quelqu'un parce qu'il t'a raconté
quelques histoires ?

— Ce n'est pas pareil.

— Si.

— Non.

Le mot lui échappa presque avant qu'il ait décidé de le prononcer.

Le visage de son père se ferma.

— Tu crois donc déjà savoir mieux que moi ce qui est bon pour le clan ?

— Je n'ai pas dit ça.

— Alors obéis.

Khaer Ghal serra les poings.

— Tu vas les tuer juste pour me punir ?

Son père s'approcha jusqu'à n'être plus qu'à quelques pas de lui.

— Non. Je vais te rappeler ce qu'ils sont.

Il désigna la sortie.

— Et demain, tu regarderas.

\scenebreak

Le lendemain, une grande partie du village se rassembla autour de la fosse.
Les habitants occupaient les aménagements disposés sur son pourtour et
discutaient en attendant l'ouverture des accès, comme ils l'avaient fait pour
des dizaines d'affrontements auparavant.

Khaer Ghal resta parmi eux.

Il avait passé une grande partie de la nuit à chercher une solution. Il avait
imaginé ouvrir l'enclos, mais les prisonniers étaient surveillés. Il avait
pensé demander de l'aide, sans parvenir à trouver qui accepterait de s'opposer
au chef pour sauver trois humains. Il avait même envisagé de retourner voir
son père avant de comprendre qu'une seconde discussion ne ferait probablement
qu'aggraver les choses.

Il ne pouvait pas ouvrir l'enclos sans être vu. Il ne pouvait pas affronter
les gardes. Il ne pouvait pas convaincre son père.

Alors il était là.

Et il détestait cela.

L'un des accès aménagés dans la paroi de la fosse s'ouvrit. Le père de Kamatsu
apparut le premier, suivi de sa mère. Kamatsu descendit derrière eux.

Khaer Ghal le repéra immédiatement.

Pendant quelques secondes, son ami chercha quelque chose parmi tous les
visages penchés au-dessus de l'arène. Son regard parcourut la foule, hésita,
puis trouva celui de Khaer Ghal.

Aucun des deux ne sourit.

Khaer Ghal aurait voulu lui expliquer qu'il avait essayé. Qu'il était allé
voir son père. Qu'il avait cherché autre chose pendant toute la nuit.

Il n'avait rien trouvé.

— Désolé...

Le mot se perdit dans le bruit autour de lui.

Kamatsu ne pouvait pas l'entendre.

L'accès se referma derrière les trois prisonniers.

De l'autre côté de la fosse, deux autres passages furent ouverts.

Les créatures apparurent.

Khaer Ghal les reconnut aussitôt. Elles appartenaient à la même espèce que
celle qu'il avait affrontée lors de sa première épreuve, quoique légèrement
plus grandes. Leur corps bas et puissant leur permettait d'accélérer
brutalement avant de bondir, tandis que leur large mâchoire pouvait infliger
en un instant des blessures suffisamment graves pour mettre fin à un combat.

Il connaissait leurs mouvements.

Kamatsu et ses parents, non.

On avait tout de même donné aux trois prisonniers de quoi se défendre. Le père
tenait une lance courte, la mère un bâton renforcé et Kamatsu une petite lame.
Entre les mains de personnes qui n'avaient jamais appris à les utiliser
ensemble, ces armes ne suffisaient pas à rétablir l'équilibre.

La famille se regroupa instinctivement.

Une créature partit vers la gauche.

L'autre contourna lentement les humains par la droite.

Khaer Ghal sentit ses doigts se refermer autour de la poignée de sa dague.

Il devait rester à sa place.

Son père le lui avait ordonné.

La première créature bondit.

Le père de Kamatsu tenta de présenter sa lance devant lui. La pointe atteignit
l'épaule de la bête, mais son geste manquait de force et d'assurance. Le choc
le fit reculer de plusieurs pas tandis que l'animal retombait déjà sur ses
pattes.

La seconde attaqua presque aussitôt.

La mère de Kamatsu poussa son fils derrière elle et abattit son bâton. La
créature esquiva le coup et recula avant de tourner de nouveau autour d'eux.

Khaer Ghal ne quittait plus leurs pattes des yeux.

Il connaissait cette manière de faire. Les bêtes testaient leurs adversaires.
Elles attaquaient, reculaient, observaient les réactions et recommençaient
jusqu'à trouver une ouverture.

Le père de Kamatsu récupéra sa lance.

La première créature se tassa.

Khaer Ghal vit son poids passer sur ses pattes arrière.

Sa tête descendit.

Il connaissait ce mouvement.

Il l'avait appris ici même.

— Bouge...

Personne ne l'entendit.

La bête bondit.

Le père de Kamatsu esquiva trop tard. Il fut heurté et projeté au sol, sa lance
roulant plusieurs pas plus loin.

Kamatsu courut vers lui.

La seconde créature changea immédiatement de direction.

Vers Kamatsu.

Khaer Ghal vit les pattes arrière se contracter.

Il sut ce qui allait arriver.

Il savait comment l'éviter.

Kamatsu, non.

Quelque chose céda.

Khaer Ghal passa la barrière.

\scenebreak

Des cris éclatèrent autour de lui.

Il n'en écouta aucun.

Il trouva l'un des accès descendant dans la fosse, franchit les dernières
marches presque en sautant, tira sa dague et courut.

— Khaer Ghal !

La voix de son père domina un instant toutes les autres.

Il continua.

— Kamatsu !

Le garçon tourna la tête.

— Derrière toi !

Kamatsu se jeta sur le côté.

La créature bondit.

Ses mâchoires claquèrent dans le vide.

Khaer Ghal arriva.

Il frappa.

Sa lame ouvrit le flanc de la bête, qui pivota aussitôt vers lui.

Cette fois, il connaissait son adversaire.

Deux ans auparavant, il avait affronté une créature semblable dans cette même
fosse. Il avait eu peur, perdu sa dague et laissé une partie de son oreille
entre ses crocs avant de comprendre comment elle attaquait.

Depuis, il s'était entraîné.

Encore et encore.

La créature se tassa.

Les pattes arrière.

La tête.

Le bond.

Khaer Ghal se décala.

La mâchoire passa devant lui.

Sa dague frappa une seconde fois.

— Restez derrière moi !

Il ne savait pas s'il s'adressait à Kamatsu ou à toute sa famille. Cela
n'avait aucune importance.

La seconde créature arrivait déjà.

Trop vite.

Khaer Ghal se jeta au sol.

Les crocs passèrent au-dessus de son épaule.

Il roula dans la poussière et se releva.

La première revenait.

Deux.

Il n'en avait jamais affronté deux.

Il recula, cherchant à conserver les deux animaux dans son champ de vision.
L'un avança. L'autre contourna.

Mauvais.

Il changea de position.

La première feinta.

La seconde bondit.

Khaer Ghal pivota trop tard.

Une patte le frappa au côté.

Il tomba sur un genou.

La mâchoire arriva.

Il leva le bras par réflexe, sentit des griffes lui ouvrir la peau et se jeta
en arrière avant que les crocs puissent se refermer.

— Khaer Ghal !

Kamatsu.

Il n'eut pas le temps de répondre.

La première créature attaquait de nouveau.

Khaer Ghal esquiva.

À peine.

Il commençait à comprendre son erreur. Tant qu'il essayait de surveiller les
deux créatures à la fois, chacune profitait du moment où son attention se
portait sur l'autre. Il ne gagnerait pas ainsi. Il devait en éliminer une
rapidement, et la première était déjà blessée.

Ce serait elle.

Khaer Ghal recula de quelques pas et la créature le suivit. Il savait que la
seconde se déplaçait quelque part sur sa droite, mais il se força à l'ignorer
pour concentrer toute son attention sur celle qui lui faisait face. Encore un
pas. La bête blessée se tassa presque aussitôt sur elle-même.

Khaer Ghal reconnut les signes : le poids sur les pattes arrière, la tête qui
s'abaissait, les muscles qui se tendaient avant le bond. Cette fois, il ne
bougea pas trop tôt. Il attendit jusqu'au dernier instant, se décala lorsque
la créature se détendit vers lui et fit remonter sa dague sous sa mâchoire.

La bête s'effondra dans la poussière.

Khaer Ghal arracha sa lame et se retourna.

Trop tard.

La seconde créature le percuta dans le dos et le projeta au sol, lui coupant
le souffle. Sa dague lui échappa des doigts et glissa hors de portée tandis
que l'animal pivotait déjà pour revenir vers lui. Khaer Ghal chercha son arme,
mais elle était trop loin.

Quelque chose tomba alors dans la poussière devant lui.

La petite lame de Kamatsu.

— Attrape !

Khaer Ghal referma aussitôt les doigts dessus. Il n'eut pas le temps de
répondre : la bête bondissait déjà. Il roula sur le côté et entendit les crocs
claquer contre la terre à l'endroit où se trouvait sa tête une fraction de
seconde plus tôt.

Il se releva et frappa presque dans le même mouvement, mais la lame glissa sur
les côtes de l'animal sans pénétrer profondément. La créature pivota et Khaer
Ghal recula, une fois, puis une seconde, retrouvant peu à peu la distance dont
il avait besoin.

La bête se tassa.

Cette fois encore, il attendit.

Sa tête descendit légèrement et Khaer Ghal reconnut le mouvement qu'il avait
déjà observé tant de fois. Il résista à l'envie de bouger, laissa la créature
s'engager entièrement dans son attaque et ne fit un pas de côté qu'au moment
où elle bondit.

Il planta la lame de toutes ses forces.

Le choc faillit lui arracher l'arme des mains. Emportée par son propre élan,
la créature s'effondra contre lui et son poids l'entraîna au sol.

Puis plus rien.

Khaer Ghal resta immobile quelques secondes, coincé sous le corps de l'animal,
le souffle court et les oreilles bourdonnantes.

Autour de lui, l'arène s'était tue.

Il repoussa la créature et se dégagea difficilement.

— Kamatsu ?

Son ami était debout. Son père et sa mère également.

Khaer Ghal expira enfin.

Puis il leva les yeux.

Son père le regardait depuis le bord de la fosse.

\scenebreak

Les règles de l'arène étaient anciennes. Un affrontement prenait fin lorsque
les adversaires désignés étaient morts ou incapables de poursuivre. Khaer Ghal
avait enfreint suffisamment de règles ce jour-là, mais son intervention n'en
changeait pas l'issue : les deux créatures étaient mortes et les trois
prisonniers étaient encore en vie.

Le combat était terminé.

Kamatsu et ses parents furent reconduits vers leur enclos sous les regards des
habitants. Khaer Ghal voulut les suivre, ne serait-ce que pour vérifier leurs
blessures et parler à son ami, mais deux guerriers l'attendaient à la sortie
de la fosse.

Son père se tenait derrière eux.

Le chef ne cria pas.

Il n'en avait pas besoin.

Khaer Ghal soutint son regard quelques secondes avant que les guerriers ne
l'emmènent vers la place centrale.

Il savait ce qui allait arriver.

Il avait désobéi à un ordre direct. Il avait interrompu un combat organisé par
le village, défié publiquement son chef et, plus grave encore aux yeux de son
père, il l'avait fait pour protéger des prisonniers après avoir reçu l'ordre
explicite de les regarder mourir.

La punition devait donc être publique elle aussi.

On l'attacha à l'un des poteaux de la place devant les habitants qui sortaient
encore de l'arène. Certains restèrent pour regarder. D'autres poursuivirent
simplement leur chemin ; un enfant puni par son père n'avait rien
d'extraordinaire dans le village, même lorsque cet enfant était le fils du
chef.

Khaer Ghal ne protesta pas.

Les premiers coups lui firent serrer les mâchoires.

Les suivants lui firent perdre le compte.

Il essaya de ne pas crier. Il ne savait même plus très bien si c'était par
fierté, par défi envers son père ou simplement parce qu'il refusait de lui
donner davantage que ce qu'il pouvait lui prendre de force.

Lorsqu'on détacha enfin ses poignets, ses jambes faillirent céder. Il resta
debout par obstination, le souffle court, le corps douloureux après le combat
et la punition.

Les habitants commencèrent à se disperser.

Son père, lui, resta et attendit que Khaer Ghal relève la tête.

— Tu crois les avoir sauvés ?

Khaer Ghal ne répondit pas.

— Tu n'as rien changé.

Son père s'approcha.

— Ils retourneront dans l'arène demain.

Khaer Ghal le fixa.

Pendant quelques secondes, il crut avoir mal entendu.

— Quoi ?

— Demain.

— Mais ils ont survécu.

— Grâce à toi.

— Ils ont survécu.

— Et ils recommenceront.

Khaer Ghal sentit revenir la colère qu'il avait éprouvée la veille.

La veille encore, il avait essayé de convaincre son père de changer sa
décision. Dans l'arène, il avait désobéi parce que Kamatsu allait mourir. Même
en sautant dans la fosse, il n'avait pensé qu'à empêcher ce qui se produisait
devant lui.

Maintenant, quelque chose était différent.

Son père pouvait recommencer demain.

Puis le jour suivant.

Il pouvait décider que Kamatsu travaillerait, combattrait ou mourrait. Il
pouvait décider de la même chose pour ses parents et pour tous ceux que les
guerriers ramenaient au village.

Simplement parce qu'ils étaient prisonniers.

Simplement parce qu'il était le chef.

Khaer Ghal repensa à une question que Kamatsu lui avait posée quelques jours
plus tôt.

\emph{Et si on ne veut pas ?}

À l'époque, il n'avait pas compris.

Maintenant, si.

Son père se pencha légèrement vers lui.

— Demain, tu resteras à ta place.

Puis il s'éloigna.

Khaer Ghal resta seul au milieu de la place, le corps douloureux et les mains
encore marquées par les liens.

Il regarda son père partir.

Puis son regard se porta vers l'enclos des prisonniers.

Demain, Kamatsu et ses parents retourneraient dans l'arène.

À moins qu'ils ne soient plus là.